\documentclass[a4paper,12pt]{article}
\usepackage[T2A]{fontenc}
\usepackage[utf8]{inputenc}
\usepackage[russian]{babel}
\usepackage{amsmath,amssymb,mathtools}
\renewcommand{\familydefault}{\sfdefault}
\usepackage{icomma}
\usepackage{geometry}
\geometry{left=23mm,right=23mm,top=21mm,bottom=22mm}
\usepackage{microtype}
\usepackage{tikz}
\usetikzlibrary{arrows.meta,positioning,shapes.geometric,calc}
\usepackage{booktabs,tabularx,array}
\usepackage{enumitem}
\usepackage{fancyhdr}
\usepackage{setspace}
\usepackage{xcolor}
\usepackage{hyperref}
\usepackage{needspace}
\definecolor{accent}{HTML}{008080}
\definecolor{darkaccent}{HTML}{006466}
\definecolor{textgray}{HTML}{373D3D}
\definecolor{muted}{HTML}{687474}
\hypersetup{colorlinks=true,linkcolor=darkaccent,urlcolor=darkaccent,
  pdftitle={Чек-лист: Доли, дроби и проценты},
  pdfauthor={Лёвин Артём Александрович}}
\setlength{\parindent}{0pt}
\setlength{\parskip}{6pt}
\onehalfspacing
\setlength{\headheight}{24pt}
\pagestyle{fancy}
\fancyhf{}
\fancyhead[L]{\footnotesize\color{muted}Доли, дроби и проценты}
\fancyhead[R]{\footnotesize\color{muted}Ксения Васильченко}
\fancyfoot[C]{\footnotesize\color{muted}\thepage}
\renewcommand{\headrulewidth}{0.4pt}
\renewcommand{\headrule}{\hbox to\headwidth{\color{accent}\leaders\hrule height \headrulewidth\hfill}}
\newcommand{\sectionline}[1]{%
  \par\Needspace{7\baselineskip}\vspace{6pt}
  {\color{darkaccent}\large\bfseries #1\par}\vspace{2pt}
  {\color{accent}\hrule height 0.55pt}\vspace{6pt}}
\newcommand{\example}[1]{\par\vspace{4pt}{\color{darkaccent}\bfseries #1}\par}
\newcommand{\important}[1]{\par{\color{darkaccent}\bfseries Главная идея.}~#1\par}
\newcommand{\task}[1]{\item #1}
\newcommand{\selfcheck}[1]{\par\noindent$\square$\enspace #1\par}
\begin{document}

\hypersetup{pageanchor=false}
\begin{titlepage}
\thispagestyle{empty}
\centering
\vspace*{29mm}
{\color{darkaccent}\large МАТЕМАТИКА\par}
\vspace{17mm}
{\Huge Чек-лист\par}
\vspace{11mm}
{\color{accent}\LARGE Доли, дроби и проценты\\[2pt]Как найти часть и целое\par}
\vspace{20mm}
{\large Математика, 6 класс\par}
\vfill
{\large Ксения Васильченко\par}
\vspace{10mm}
{11 ноября 2026 г.\par}
\vspace{18mm}
{\footnotesize Лёвин Артём Александрович\par}
{\footnotesize эксклюзивно для Ксении Васильченко\par}
\vspace{16mm}
\begin{tikzpicture}
  \draw[accent,line width=.9pt] (-5,0)
    .. controls (-3.7,.65) and (-1.7,-.60) .. (0,0)
    .. controls (1.65,.65) and (3.6,-.60) .. (5,0);
\end{tikzpicture}
\vspace{1mm}
\end{titlepage}
\hypersetup{pageanchor=true}

\sectionline{1. Целое, доля и часть}
Всю величину (книгу, массу, путь, цену) принимаем за \textbf{единицу}. Если целое разбито на равные доли, дробь $\frac{a}{b}$ показывает, что взято $a$ таких долей из $b$.

\begin{center}
\renewcommand{\arraystretch}{1.55}
\begin{tabularx}{.96\linewidth}{@{}>{\raggedright\arraybackslash}X>{\raggedright\arraybackslash}X@{}}
\toprule
Что дано & Что делаем \\
\midrule
Целое и доля & \textbf{Часть} $=$ целое $\cdot$ доля \\
Часть и её доля & \textbf{Целое} $=$ часть $:$ доля \\
\bottomrule
\end{tabularx}
\end{center}

\example{Пример 1. Известно целое}
Найдём $\frac58$ от $48$ страниц:
\[
48\cdot\frac58=\frac{48}{8}\cdot5=6\cdot5=30\ \text{страниц}.
\]
\example{Пример 2. Известна часть}
$21$ страница составляет $\frac37$ книги. Найдём объём книги:
\[
21:\frac37=21\cdot\frac73=7\cdot7=49\ \text{страниц}.
\]
\important{Перед вычислением установите, какое число обозначает всё целое, а какое --- его часть. При поиске целого делитель-доля должен отличаться от нуля.}

\sectionline{2. Проценты --- сотые доли}
Один процент равен одной сотой: $1\%=\frac1{100}$. Поэтому $p\%$ соответствует дроби $\frac p{100}$.
\[
15\%=\frac{15}{100}=\frac3{20},\qquad
25\%=\frac14,\qquad
0{,}4=\frac{40}{100}=40\%.
\]
\example{Пример 3. Найти процент от числа}
$25\%$ от $72$ рублей:
\[
72\cdot\frac{25}{100}=72\cdot\frac14=18\ \text{руб.}
\]
\example{Пример 4. Найти всё число по его проценту}
$42$ килограмма составляют $35\%$ груза. Тогда масса всего груза:
\[
42:\frac{35}{100}=42:\frac7{20}
=42\cdot\frac{20}{7}=120\ \text{кг}.
\]

\sectionline{3. Удобная запись дробей перед вычислением}
Смешанное число сначала переводят в неправильную дробь. Десятичную дробь тоже можно записать как обыкновенную:
\[
1\frac14=\frac{1\cdot4+1}{4}=\frac54,\qquad
0{,}75=\frac{75}{100}=\frac34.
\]
\example{Пример 5. Выражение со смешанным числом}
Найдём $1\frac12$ от числа $36$:
\[
36\cdot1\frac12=36\cdot\frac32=54.
\]
Это $150\%$ исходного числа. Дробь больше единицы описывает величину, превосходящую исходную.

\sectionline{4. Переводим текст задачи на язык дробей}
Формулировки «половина», «четверть», «треть» означают соответственно $\frac12$, $\frac14$, $\frac13$.

\example{Пример 6. Несколько частей одного целого}
Прочитали четверть книги и ещё три восьмых \emph{всей} книги. Всего --- $50$ страниц. Тогда прочитанная доля:
\[
\frac14+\frac38=\frac28+\frac38=\frac58.
\]
Число страниц в книге:
\[
50:\frac58=50\cdot\frac85=80\ \text{страниц}.
\]
\important{Складывать доли сразу можно, если они взяты от одного и того же целого и описывают непересекающиеся части.}

\sectionline{5. Остаток: отнимаем долю от единицы}
Если израсходовали $\frac27$ всего запаса, осталось:
\[
1-\frac27=\frac77-\frac27=\frac57.
\]
\begin{center}
\begin{tikzpicture}[x=1.45cm,y=1cm]
\fill[accent!17] (0,0) rectangle (2,.42);
\fill[muted!12] (2,0) rectangle (7,.42);
\draw[accent,line width=.85pt] (0,0) rectangle (7,.42);
\foreach \x in {1,...,6}{\draw[accent!65,line width=.55pt] (\x,0)--(\x,.42);}
\node[anchor=south] at (1,.50) {$\frac27$ израсходовано};
\node[anchor=south] at (4.5,.50) {$\frac57$ осталось};
\node[anchor=north] at (3.5,-.08) {Весь запас --- $1$};
\end{tikzpicture}
\end{center}
\example{Пример 7. Найти первоначальную массу}
После продажи $\frac27$ картофеля осталось $25$ кг. Эти $25$ кг составляют $\frac57$ всей массы, поэтому
\[
25:\frac57=25\cdot\frac75=35\ \text{кг}.
\]

\sectionline{6. Важное различие: доля от всего и доля от остатка}
Слова \textbf{«четверть оставшегося»} означают четверть нового, уменьшившегося количества.
\example{Пример 8. Два последовательных действия}
Было $120$ кг яблок. Сначала продали $\frac13$ всего запаса, затем $\frac14$ \emph{оставшегося}. Найдём остаток.

После первой продажи осталось:
\[
120\cdot\left(1-\frac13\right)=120\cdot\frac23=80\ \text{кг}.
\]
После второй продажи осталось:
\[
80\cdot\left(1-\frac14\right)=80\cdot\frac34=60\ \text{кг}.
\]
\important{Во втором действии четверть берётся от $80$ кг. Выражение $1-\frac13-\frac14$ здесь неверно: доли относятся к разным количествам.}

\sectionline{7. Скидка: уменьшение и новая цена}
При снижении цены на $r\%$ новая цена составляет $(100-r)\%$ исходной.
\example{Пример 9. Известен размер скидки в рублях}
Цена уменьшилась на $20\%$, скидка составила $30$ рублей. Значит, первоначальная цена:
\[
30:\frac{20}{100}=30:\frac15=150\ \text{руб.}
\]
После скидки товар стоит $150-30=120$ рублей.

\example{Пример 10. Известна цена после скидки}
После скидки $20\%$ товар стоит $96$ рублей. Новая цена составляет $80\%=\frac45$ первоначальной, поэтому
\[
96:\frac45=96\cdot\frac54=120\ \text{руб.}
\]
Проверка: $20\%$ от $120$ равны $24$; $120-24=96$.

\par\medskip
{\color{darkaccent}\bfseries Типичные ошибки.} При поиске целого умножают вместо деления; путают скидку с новой ценой; вычитают доли от разных количеств.

\clearpage
\thispagestyle{fancy}
\begin{center}
{\color{darkaccent}\Large\bfseries Алгоритм: как найти часть или целое}\par
\vspace{9mm}
\end{center}
\begin{center}
\begin{tikzpicture}[
  >=Latex,
  action/.style={draw=accent,rounded corners=2mm,align=center,text width=11.2cm,minimum height=1.15cm,inner sep=5pt},
  check/.style={diamond,draw=darkaccent,aspect=2.6,align=center,text width=3.1cm,inner sep=3pt},
  branch/.style={draw=accent,rounded corners=2mm,align=center,text width=5.3cm,minimum height=1.4cm,inner sep=4pt},
  end/.style={ellipse,draw=darkaccent,align=center,minimum width=5.8cm,minimum height=.85cm},
  arrow/.style={->,draw=textgray,thick},
  every node/.style={font=\normalsize}
]
\node[end] (a) at (0,0) {Прочитайте условие};
\node[action] (b) at (0,-1.95) {Примите всё количество за $1$. Переведите проценты и слова в дроби};
\node[action] (c) at (0,-4.15) {Определите долю той части, для которой дано число};
\node[check] (d) at (0,-6.85) {Известно целое?};
\node[branch] (e) at (-3.35,-9.65) {Найдите часть:\\ целое $\cdot$ доля};
\node[branch] (f) at (3.35,-9.65) {Найдите целое:\\ часть $:$ доля};
\node[action] (g) at (0,-12.2) {Проверьте вычисление, единицы и смысл ответа};
\node[end] (h) at (0,-14.15) {Запишите ответ};
\draw[arrow] (a)--(b);
\draw[arrow] (b)--(c);
\draw[arrow] (c)--(d);
\draw[arrow] (d.west)-|node[pos=.26,above]{да}(e.north);
\draw[arrow] (d.east)-|node[pos=.26,above]{нет}(f.north);
\draw[arrow] (e.south)--++(0,-.38)-|(g.west);
\draw[arrow] (f.south)--++(0,-.38)-|(g.east);
\draw[arrow] (g)--(h);
\end{tikzpicture}
\end{center}

\clearpage
\sectionline{Тренировочный блок \hfill Путь А}
\textcolor{muted}{10 заданий: от простых вычислений к многошаговым задачам. В текстовых задачах записывайте, какую долю составляет известная величина.}
\vspace{4mm}
\begin{enumerate}[label=\textbf{\arabic*.},leftmargin=8mm,itemsep=13pt]
\task{Найдите $\frac38$ от числа $72$.}
\task{Число $24$ составляет $\frac25$ некоторого числа. Найдите это число.}
\task{Вычислите $1\frac12$ от числа $36$.}
\task{Найдите $18\%$ от числа $250$.}
\task{$45$ страниц составляют $15\%$ книги. Сколько страниц в книге?}
\task{В книге $90$ страниц. Прочитали $\frac25$ книги. Сколько страниц осталось прочитать?}
\task{Прочитали четверть книги и ещё три восьмых всей книги --- всего $75$ страниц. Сколько страниц в книге?}
\task{В магазине было $84$ кг ягод. Сначала продали $\frac13$ всех ягод, затем $\frac14$ оставшихся. Сколько килограммов ягод осталось?}
\task{После скидки $15\%$ товар стоит $340$ рублей. Какова была его первоначальная цена?}
\task{В первый день турист прошёл $35\%$ всего пути, во второй --- $\frac25$ пути, оставшегося после первого дня. После двух дней осталось пройти $78$ км. Какова длина всего пути?}
\end{enumerate}

\clearpage
\sectionline{Ответы к тренировочному блоку}
\vspace{4mm}
\begin{center}
\renewcommand{\arraystretch}{1.7}
\begin{tabularx}{.90\linewidth}{@{}>{\centering\arraybackslash}p{18mm}X@{}}
\toprule
\textbf{№} & \textbf{Ответ}\\
\midrule
1 & $27$\\
2 & $60$\\
3 & $54$\\
4 & $45$\\
5 & $300$ страниц\\
6 & $54$ страницы\\
7 & $120$ страниц\\
8 & $42$ кг\\
9 & $400$ рублей\\
10 & $200$ км\\
\bottomrule
\end{tabularx}
\end{center}
\vspace{14mm}
{\color{darkaccent}\large\bfseries Итоговая самопроверка}\par\vspace{5mm}
\selfcheck{Могу найти дробь или процент от известного целого.}
\selfcheck{Могу восстановить целое по известной части.}
\selfcheck{Умею переводить смешанные и десятичные числа в дроби.}
\selfcheck{Отличаю долю от целого от доли оставшегося.}
\selfcheck{Различаю размер скидки и цену после скидки.}
\selfcheck{Могу проверить, правдоподобен ли полученный ответ.}
\vfill
\begin{center}
{\footnotesize\color{muted}Лёвин Артём Александрович --- эксклюзивно для Ксении Васильченко}
\end{center}
\end{document}
